\twocolumn[
    {\centering
        {\LARGE\bfseries Spatial Distribution of Bidirectional Charging Stations by Simulation and Evaluation of Traffic and Grid Data\par}
        \vspace{1.2em}

        {\large
        Achim\,\textsc{Baumgärtner}\textsuperscript{1},\quad
        Erik\,\textsc{Weigand}\textsuperscript{1},\quad \\
        Gian Bruno\,\textsc{Truöl}\textsuperscript{1,2},\quad
        Mirko\,\textsc{Sonntag}\textsuperscript{1,2}\par
        }
        \vspace{0.6em}

        {\small
        \textsuperscript{1}University of Applied Sciences Esslingen,
            Esslingen University, Germany\par
        \textsuperscript{2}Fraunhofer Application Center Keim,
            Esslingen, Germany\par
        }
        \vspace{0.4em}

        {\small\texttt{acbait01@hs-esslingen.de, erwewi01@hs-esslingen.de, \\ gian-bruno.truoel@hs-esslingen.de, mirko.sonntag@hs-esslingen.de}\par}
        \vspace{0.8em}

        {\small Submitted: 01 August 2026 \quad|\quad
            Revised: 25 March 2026 \quad|\quad
            Accepted: 01 August 2026\par}
    }

    \vspace{0.8em}\hrule\vspace{0.8em}

    % =========================================
    % ABSTRACT
    % =========================================
    DIESER AUSZUG DES JOURNAL ARTIKELS STELLT NUR EINEN DRAFT DAR.


    {\noindent\textbf{Abstract.} This paper investigates how bidirectional charging stations can be distributed across an urban area by combining traffic simulation with power-grid indicators. We model charging demand as a spatiotemporal process and evaluate candidate locations with respect to accessibility, waiting time, and local grid stress. The proposed workflow integrates mobility traces, station utilization forecasts, and feeder-level constraints to identify robust deployment strategies. Results indicate that jointly optimizing traffic and grid criteria can reduce peak congestion while improving charging availability in high-demand districts.\par}

    \vspace{0.4em}
    \noindent\textbf{Keywords:}
    Bidirectional charging, traffic simulation, real-time modeling,
    neural networks, Python

    \vspace{0.8em}\hrule\vspace{1.2em}
]

% =============================================
% INHALTSVERZEICHNIS
% Hinweis: In Journal-Papers oft weggelassen.
% For a technical report / extended paper,
% this can still be useful.
% =============================================
% \tableofcontents
% \newpage
